% ==================================================================
%  PART 3 --- continues Chapter 3 (Section: Core Algorithms onward),
%  then Chapter 4 (Complexity Analysis) and Chapter 5 (Worked Example)
%  Paste this directly after the last line of part2_ch1-3.tex
%  ("...must preserve on exactly the nodes where it could have been
%  broken.") and before whatever chapter currently follows it.
% ==================================================================

\section{Core Algorithms}
\label{sec:core-algorithms}
Every operation on a Segment Tree is a recursive walk over the same implicit tree described in Section~\ref{sec:core-algorithms}: at each node the algorithm asks a cheap question about the current interval $[lo,hi]$ and either answers immediately, recurses into one child, or recurses into both. This section develops the three core routines --- build, update, and query --- one at a time, and Section~\ref{sec:query-example} traces each of them through the running $n=8$ example.

\subsection{Build}
\label{sec:build}
Construction populates every node's stored summary, starting from the leaves and merging upward. Algorithm~\ref{alg:build} implements it as a single recursive pass.

\begin{lstlisting}[style=cpp, caption={Recursive build.}, label={alg:build}]
void build(int node, int lo, int hi) {
    if (lo == hi) {
        tree[node] = A[lo];
        return;
    }
    int mid = (lo + hi) / 2;
    build(2*node,   lo,     mid);
    build(2*node+1, mid+1,  hi);
    tree[node] = tree[2*node] + tree[2*node+1];   // merge
}
\end{lstlisting}

\begin{keyidea}
Build never inspects a query range at all --- it simply materializes the invariant of Equation~\eqref{eq:invariant} everywhere, once, so that later queries can trust every node's stored value without recomputation.
\end{keyidea}

Because each recursive call either reaches a leaf (processing one array element) or splits into exactly two children, the call tree of \texttt{build} is precisely the Segment Tree itself: one call per node, $2n-1$ calls in total by Lemma~\ref{lem:nodes}. Each call does $O(1)$ work beyond its two recursive calls, so:

\begin{lem}[Build time]\label{lem:build-time}
\texttt{build} runs in $\Oh{n}$ time and uses $\Oh{n}$ auxiliary space for the \texttt{tree} array.
\end{lem}
\begin{proof}
The number of calls is $2n-1$ (Lemma~\ref{lem:nodes}), and each call performs $O(1)$ work outside of its recursive calls, so the total work is $\Oh{2n-1}=\Oh{n}$. The recursion depth is the tree height $\Oh{\log n}$ (Lemma~\ref{lem:height}), which bounds the auxiliary stack space, while the \texttt{tree} array itself occupies $\Oh{n}$ cells.
\end{proof}

\subsection{Point Update}
\label{sec:update}
A point update changes $A[pos]$ to a new value and must restore the invariant of Equation~\eqref{eq:invariant} at every node whose interval contains $pos$. Crucially, that set of nodes is not arbitrary: it is exactly the nodes on the root-to-leaf path for $pos$.

\begin{lem}[Update touches one path]\label{lem:update-path}
For a fixed index $pos$, the set of nodes whose interval contains $pos$ forms a single path from the root to the leaf $[pos,pos]$.
\end{lem}
\begin{proof}
By Lemma~\ref{lem:partition}, the intervals at any one level are pairwise disjoint, so at most one node per level can contain $pos$. The root's interval $[1,n]$ always contains $pos$, and if a node at depth $d$ contains $pos$ then exactly one of its two children's intervals (by construction, a partition of the parent's interval) contains $pos$ as well, unless the node is already the leaf $[pos,pos]$. Hence there is exactly one such node at every depth from $0$ to $\lceil\log_2 n\rceil$, and consecutive ones are parent and child, i.e.\ a path.
\end{proof}

This lemma is the entire justification for the update algorithm: descend toward the leaf containing $pos$, ignoring the sibling subtree at every step, then recompute merges on the way back up.

\begin{lstlisting}[style=cpp, caption={Point update.}, label={alg:update}]
void update(int node, int lo, int hi, int pos, int val) {
    if (lo == hi) {                 // pos == lo == hi
        tree[node] = val;
        return;
    }
    int mid = (lo + hi) / 2;
    if (pos <= mid)
        update(2*node,   lo,    mid, pos, val);
    else
        update(2*node+1, mid+1, hi,  pos, val);
    tree[node] = tree[2*node] + tree[2*node+1];   // re-merge on the way back up
}
\end{lstlisting}

\begin{pitfall}
It is tempting to update only the leaf and stop. This leaves every ancestor's stored summary stale, silently corrupting every future query that touches $pos$. The re-merge line at the end of \texttt{update} is not optional bookkeeping --- it is what restores Equation~\eqref{eq:invariant}.
\end{pitfall}

\begin{lem}[Update time]\label{lem:update-time}
\texttt{update} runs in $\Oh{\log n}$ time.
\end{lem}
\begin{proof}
By Lemma~\ref{lem:update-path}, the recursion visits exactly one node per level, and the tree has $\lceil\log_2 n\rceil+1$ levels by Lemma~\ref{lem:height}. Each visited node performs $O(1)$ work (a comparison, a recursive call, and a merge), so the total work is $\Oh{\log n}$.
\end{proof}

\subsection{Range Query}
\label{sec:query}
A range query for $[ql,qr]$ should return $S(ql,qr) = \sum_{i=ql}^{qr} A[i]$ without touching every element. The algorithm compares the current node's interval $[lo,hi]$ against $[ql,qr]$ and falls into exactly one of three cases at each step.

\begin{table}[h]
\centering
\caption{The three cases handled at each recursive step of a range query.}
\label{tab:query-cases}
\begin{tabular}{>{\raggedright\arraybackslash}p{3.0cm}>{\raggedright\arraybackslash}p{4.6cm}p{5.0cm}}
\toprule
\tblhead{Case} & \tblhead{Condition} & \tblhead{Action} \\
\midrule
No overlap      & $qr < lo$ or $hi < ql$            & Return the identity; contributes nothing. \\
Total overlap   & $ql \le lo$ and $hi \le qr$        & Return \texttt{tree[node]} directly; no need to look deeper. \\
Partial overlap & otherwise                          & Recurse into both children and merge their results. \\
\bottomrule
\end{tabular}
\end{table}

\begin{lstlisting}[style=cpp, caption={Range query.}, label={alg:query}]
int query(int node, int lo, int hi, int ql, int qr) {
    if (qr < lo || hi < ql)
        return 0;                                  // no overlap (identity for sum)
    if (ql <= lo && hi <= qr)
        return tree[node];                          // total overlap
    int mid = (lo + hi) / 2;
    int left  = query(2*node,   lo,    mid, ql, qr);
    int right = query(2*node+1, mid+1, hi,  ql, qr);
    return left + right;                             // partial overlap
}
\end{lstlisting}

\begin{keyidea}
The identity element (here, $0$ for sum; $+\infty$ for minimum, $-\infty$ for maximum, $0$ for GCD) is what lets the no-overlap case return \emph{something} without a special-cased branch at the call site: merging with the identity never changes the result.
\end{keyidea}

Unlike \texttt{build} and \texttt{update}, a query does not visit a bounded number of nodes per level in general --- it can visit up to four per level, because a query range can cut into an interval from both sides. The next subsection makes this precise, since it is the crux of the $\Oh{\log n}$ query bound and easy to get wrong by assuming the recursion always halves the range.

\subsection{Why Range Query Is \texorpdfstring{$O(\log n)$}{O(log n)}: the Overlap Argument}
\label{sec:overlap-argument}
Call a node \emph{fully covered} if it falls into the total-overlap case of Table~\ref{tab:query-cases}, and \emph{boundary} if it falls into the partial-overlap case (its interval straddles $ql$ or $qr$). No-overlap nodes are pruned immediately and never recurse further, so they contribute no cost beyond the $O(1)$ check that discards them.

\begin{thm}[Query time]\label{thm:query-time}
\texttt{query} performs $\Oh{\log n}$ merge operations and visits $\Oh{\log n}$ nodes in total.
\end{thm}
\begin{proof}
At any fixed level $d$ of the tree, at most two nodes can be boundary nodes: one containing the left endpoint $ql$ and one containing the right endpoint $qr$ (there is at most one node per level containing any single index, by the disjointness half of Lemma~\ref{lem:partition}, and a node cannot straddle a query boundary unless it contains $ql$ or $qr$). A boundary node is exactly the kind of node that recurses into both children, so the number of boundary nodes at level $d+1$ descended from level $d$'s boundary nodes is at most $2\cdot 2 = 4$, but by the same argument only at most $2$ of those can again be boundary at level $d+1$ --- the rest are resolved as fully-covered or no-overlap immediately, which do not recurse further. Formally: define $B(d)$ as the number of boundary nodes at level $d$. Then $B(0)\le 1$ and $B(d)\le 2$ for all $d\ge 1$, because a boundary node has at most two children and at most one of a node's two children can itself be a boundary node on each side --- so the left boundary chain and the right boundary chain each advance by at most one boundary node per level. Since a fully-covered node returns immediately without recursing, the total number of nodes visited at any level is $\Oh{1}$, and the recursion depth is $\Oh{\log n}$ by Lemma~\ref{lem:height}. Summing $\Oh{1}$ work over $\Oh{\log n}$ levels gives $\Oh{\log n}$ total.
\end{proof}

\begin{remark}
An equivalent, more concrete way to see Theorem~\ref{thm:query-time}: any interval $[ql,qr]$ decomposes into at most $2\lceil \log_2 n\rceil$ dyadic node-intervals, since the decomposition can be built by walking $ql$ and $qr$ upward simultaneously, each accumulating at most one fully-covered node per level. This is the same bound used informally in the Key Idea box of Section~\ref{ch:segtree}.
\end{remark}

\begin{figure}[H]
\centering
\begin{tikzpicture}[node distance=7mm and 6mm]
\node[cutnode] (root) {$[1,8]$};
\node[pathnode, below left=of root] (l) {$[1,4]$};
\node[hitnode, below right=of root] (r) {$[5,8]$};
\node[hitleaf, below left=of l] (a) {$[1,2]$};
\node[pathleaf, below right=of l] (b) {$[3,4]$};
\node[cutleaf, below left=of r] (c) {$[5,6]$};
\node[cutleaf, below right=of r] (d) {$[7,8]$};
\draw[segedge] (root)--(l) (root)--(r) (l)--(a) (l)--(b) (r)--(c) (r)--(d);
\end{tikzpicture}
\caption{Query $[1,4]$ against the array of Figure~\ref{fig:hierarchy}. Green nodes are fully covered and returned directly; gold nodes are boundary nodes that recurse further; grey dashed nodes are pruned by the no-overlap case and never visited beyond the check itself. Only $[1,2]$ and $[3,4]$ are actually summed.}
\label{fig:query-example}
\end{figure}

% ==================================================================
\chapter{Complexity Analysis}
\label{ch:complexity}

\section{Summary of Results}
Chapter~\ref{ch:segtree} established the three core algorithms and proved a time bound for each individually. This short chapter collects those results, states the corresponding space bound, and discusses the constants that Table~\ref{tab:heights} already hinted at.

\begin{table}[h]
\centering
\caption{Time and space complexity of the core Segment Tree operations, as proved in Chapter~\ref{ch:segtree}.}
\label{tab:complexity-ch4}
\begin{tabular}{llp{6.3cm}}
\toprule
\tblhead{Operation} & \tblhead{Time} & \tblhead{Source} \\
\midrule
Build          & $\Oh{n}$        & Lemma~\ref{lem:build-time} \\
Point update   & $\Oh{\log n}$   & Lemma~\ref{lem:update-time} \\
Range query    & $\Oh{\log n}$   & Theorem~\ref{thm:query-time} \\
\bottomrule
\end{tabular}
\end{table}

\section{Space Complexity}
The \texttt{tree} array is indexed with the heap-style scheme of Equation~\eqref{eq:heap}, so its size must accommodate the largest index any node can receive rather than merely the true node count of $2n-1$ (Lemma~\ref{lem:nodes}).

\begin{lem}[Space bound]\label{lem:space}
Allocating $\Oh{n}$ cells --- specifically $4n$, per Table~\ref{tab:heights} --- suffices for the \texttt{tree} array regardless of whether $n$ is a power of two.
\end{lem}
\begin{proof}
By Lemma~\ref{lem:height} the tree has height $h=\lceil\log_2 n\rceil$, so every node index under the numbering of Equation~\eqref{eq:heap} is strictly below $2^{h+1}$. Since $2^{h}<2n$ (because $h=\lceil\log_2 n\rceil$ implies $2^{h-1}<n\le 2^h$), it follows that $2^{h+1}<4n$, so $4n$ cells is a safe --- if slightly generous --- upper bound.
\end{proof}

\section{Why the Bounds Do Not Depend on the Query Pattern}
A recurring point of confusion is whether an adversarial sequence of queries and updates could push the cost above $\Oh{\log n}$ per operation. It cannot: Lemma~\ref{lem:update-path} shows the update path is determined purely by $pos$, and Theorem~\ref{thm:query-time}'s boundary-node argument bounds the number of visited nodes purely by the tree's height, independent of which values are stored or which merge results are equal to one another. The bounds are therefore \emph{worst-case per operation}, not amortized, which is one of the practical advantages of a Segment Tree over structures that only offer amortized guarantees.

\begin{remark}
This worst-case, per-operation guarantee is also why Segment Trees are preferred in settings with hard latency requirements, such as the financial-instrument scenario of Section~\ref{sec:motivation}: an amortized bound could in principle let a single unlucky query take $\Oh{n}$ time, which a Segment Tree never does.
\end{remark}

% ==================================================================
\chapter{Worked Example}
\label{ch:worked-example}

This chapter traces build, several updates, and several queries through one concrete array, so that the abstract algorithms of Chapter~\ref{ch:segtree} can be checked step by step against explicit numbers.

\section{Setup}
\label{sec:query-example}
Let $n=8$ and
\[
A = [\,3,\ 1,\ 4,\ 1,\ 5,\ 9,\ 2,\ 6\,],
\]
the same array underlying Figure~\ref{fig:hierarchy} and Figure~\ref{fig:layout}, indexed $A[1],\dots,A[8]$.

\subsection{Build Trace}
Table~\ref{tab:build-trace} lists the stored sum at every node after \texttt{build(1,1,8)} completes, using the heap-style indices of Equation~\eqref{eq:heap}.

\begin{table}[h]
\centering
\caption{Node contents after \texttt{build} on $A=[3,1,4,1,5,9,2,6]$.}
\label{tab:build-trace}
\begin{tabular}{cccc}
\toprule
\tblhead{Node} & \tblhead{Interval} & \tblhead{Sum} & \tblhead{Level} \\
\midrule
1  & $[1,8]$ & 31 & 0 \\
2  & $[1,4]$ & 9  & 1 \\
3  & $[5,8]$ & 22 & 1 \\
4  & $[1,2]$ & 4  & 2 \\
5  & $[3,4]$ & 5  & 2 \\
6  & $[5,6]$ & 14 & 2 \\
7  & $[7,8]$ & 8  & 2 \\
8  & $[1,1]$ & 3  & 3 \\
9  & $[2,2]$ & 1  & 3 \\
10 & $[3,3]$ & 4  & 3 \\
11 & $[4,4]$ & 1  & 3 \\
12 & $[5,5]$ & 5  & 3 \\
13 & $[6,6]$ & 9  & 3 \\
14 & $[7,7]$ & 2  & 3 \\
15 & $[8,8]$ & 6  & 3 \\
\bottomrule
\end{tabular}
\end{table}

Each leaf (nodes 8--15) is filled directly from $A$; each internal node is then the sum of its two children, matching Equation~\eqref{eq:merge}. The root value $31$ is the total of all eight elements, as expected.

\section{A Range Query}
Consider $\Qry(3,7)$. Table~\ref{tab:query-trace} follows the recursion of Algorithm~\ref{alg:query} call by call.

\begin{table}[h]
\centering
\caption{Recursion trace of \texttt{query(1,1,8,3,7)}.}
\label{tab:query-trace}
\begin{tabular}{cccp{5.0cm}}
\toprule
\tblhead{Node} & \tblhead{Interval} & \tblhead{Case} & \tblhead{Action} \\
\midrule
1 & $[1,8]$ & Partial & Recurse into nodes 2 and 3. \\
2 & $[1,4]$ & Partial & Recurse into nodes 4 and 5. \\
3 & $[5,8]$ & Partial & Recurse into nodes 6 and 7. \\
4 & $[1,2]$ & No overlap ($hi=2<ql=3$) & Return $0$. \\
5 & $[3,4]$ & Total overlap & Return $\text{tree}[5]=5$. \\
6 & $[5,6]$ & Total overlap & Return $\text{tree}[6]=14$. \\
7 & $[7,8]$ & Partial & Recurse into nodes 14 and 15. \\
14 & $[7,7]$ & Total overlap & Return $\text{tree}[14]=2$. \\
15 & $[8,8]$ & No overlap ($lo=8>qr=7$) & Return $0$. \\
\bottomrule
\end{tabular}
\end{table}

Combining the returned values bottom-up: node 7 returns $2+0=2$; node 3 returns $14+2=16$; node 2 returns $0+5=5$; node 1 returns $5+16=21$. Indeed $A[3]+A[4]+A[5]+A[6]+A[7]=4+1+5+9+2=21$, confirming the trace. Only $9$ of the $15$ nodes were visited, and only two nodes ($5$ and $6$) contributed a value directly without further recursion --- exactly the fully-covered decomposition described in the remark after Theorem~\ref{thm:query-time}.

\section{A Point Update}
Now apply $\Upd(6,7)$, changing $A[6]$ from $9$ to $7$. By Lemma~\ref{lem:update-path} this touches exactly the path from the root to leaf $[6,6]$, i.e.\ nodes $1\to 3\to 6\to 13$.

\begin{table}[h]
\centering
\caption{Recursion trace of \texttt{update(1,1,8,6,7)}; only the path to leaf 13 is touched.}
\label{tab:update-trace}
\begin{tabular}{ccp{8.4cm}}
\toprule
\tblhead{Node} & \tblhead{Interval} & \tblhead{Effect} \\
\midrule
13 & $[6,6]$ & Base case; $\text{tree}[13]\gets 7$. \\
6  & $[5,6]$ & Re-merge: $\text{tree}[6]\gets \text{tree}[12]+\text{tree}[13]=5+7=12$. \\
3  & $[5,8]$ & Re-merge: $\text{tree}[3]\gets \text{tree}[6]+\text{tree}[7]=12+8=20$. \\
1  & $[1,8]$ & Re-merge: $\text{tree}[1]\gets \text{tree}[2]+\text{tree}[3]=9+20=29$. \\
\bottomrule
\end{tabular}
\end{table}

Nodes $2,4,5,7$ and every leaf other than $13$ are untouched, since none of their intervals contain index $6$ --- exactly as Lemma~\ref{lem:update-path} predicts. Re-running $\Qry(3,7)$ after this update would now correctly return $4+1+5+7+2=19$, reflecting the updated value without any node outside the single path having been revisited.

\begin{remark}
Table~\ref{tab:update-trace} and Table~\ref{tab:query-trace} together illustrate the asymmetry between the two operations: an update always touches a fixed-size path of $\Oh{\log n}$ nodes regardless of which index changed, while a query's cost depends on how the range $[ql,qr]$ happens to align with the dyadic intervals of the tree --- though it is $\Oh{\log n}$ in every case, by Theorem~\ref{thm:query-time}.
\end{remark}
